% Part: first-order-logic
% Chapter: completeness
% Section: lindenbaums-lemma

\documentclass[../../../include/open-logic-section]{subfiles}

\begin{document}

\iftag{FOL}
      {\olfileid[zh]{fol}{com}{lin}}
      {\olfileid[zh]{pl}{com}{lin}}

\olsection{Lindenbaum 引理}

\begin{explain}
我们现在证明一个引理，它表明任一致的!!{sentence}集都包含于某个不仅是
一致的、而且也是!!{complete}的语句集。该证明通过一次添加一个!!{sentence}来进行，并保证在每一步骤中所得集合仍然一致。我们这样做是为了
对每个 $!A$，要么 $!A$ 要么 $\lnot !A$ 在某个阶段被添加进来。该构造中所有
阶段的并集于是包含 $!A$ 或其否定~$\lnot !A$，因而是完全的。它也是一致的，
因为我们确保在每个阶段都不引入不一致性。
\end{explain}

\begin{lem}[Lindenbaum 引理]
\ollabel{lem:lindenbaum} 语言~$\Lang{L}$ 中的每个一致
集合~$\Gamma$ 都可被扩张为!!a{complete}且一致的集合~$\Gamma^*$。
\end{lem}

\begin{proof}
设 $\Gamma$ 一致。设 $!A_0$、$!A_1$、
\dots{} 是~$\Lang L$ 的所有!!{sentence}的一个枚举。
定义 $\Gamma_0 = \Gamma$，且
\[
\Gamma_{n+1} =
\begin{cases}
\Gamma_n \cup \{ !A_n \} & \textrm{若 $\Gamma_n \cup \{!A_n\}$
  一致；} \\
\Gamma_n \cup \{ \lnot !A_n \} & \textrm{否则。}
\end{cases}
\]
设 $\Gamma^* = \bigcup_{n \geq 0} \Gamma_n$。

每个 $\Gamma_n$ 都是一致的：$\Gamma_0$ 按定义一致。
若 $\Gamma_{n+1} = \Gamma_n \cup \{!A_n\}$，这是因为后者
一致。若它不一致，则 $\Gamma_{n+1} = \Gamma_n \cup \{\lnot
!A_n\}$。我们必须验证 $\Gamma_n \cup \{\lnot !A_n\}$ 一致。设它不一致。
那么 \emph{二者} $\Gamma_n \cup
\{!A_n\}$ 与 $\Gamma_n \cup \{\lnot !A_n\}$ 都不一致。这意味着 $\Gamma_n$ 将
因
\iftag{FOL}{%
  \tagrefs{prfAX/{fol:axd:prv:prop:provability-exhaustive},
    prfSC/{fol:seq:prv:prop:provability-exhaustive},
    prfND/{fol:ntd:prv:prop:provability-exhaustive},
    prfTab/{fol:tab:prv:prop:provability-exhaustive}}}{%
  \tagrefs{prfAX/{pl:axd:prv:prop:provability-exhaustive},
    prfSC/{pl:seq:prv:prop:provability-exhaustive},
    prfND/{pl:ntd:prv:prop:provability-exhaustive},
    prfTab/{pl:tab:prv:prop:provability-exhaustive}}}而
不一致，这与归纳假设相反。

对每个~$n$ 及每个 $i < n$，$\Gamma_i \subseteq \Gamma_n$。这
通过对~$n$ 的简单归纳得出。对于 $n=0$，不存在 $i < 0$，
因此该结论自动成立。对于归纳步，设它对~$n$ 成立。我们证明若 $i < n+1$ 则
$\Gamma_i \subseteq \Gamma_{n+1}$。由构造有 $\Gamma_{n+1} = \Gamma_n \cup
\{!A_n\}$ 或 $= \Gamma_n \cup \{\lnot !A_n\}$。故 $\Gamma_n \subseteq
\Gamma_{n+1}$。若 $i < n+1$，则由归纳假设（若 $i < n$）或显然的事实
$\Gamma_n \subseteq \Gamma_n$（若 $i = n$）有 $\Gamma_i \subseteq
\Gamma_n$。于是由~$\subseteq$ 的传递性得 $\Gamma_i \subseteq
\Gamma_{n+1}$。

由此可得 $\Gamma^*$ 一致。理由如下：设 $\Gamma' \subseteq \Gamma^*$
有限。每个 $!B \in \Gamma'$ 也都在某个~$\Gamma_i$（对某个~$i$）中。设 $n$
为其中最大的。由于若 $i \le n$ 则 $\Gamma_i \subseteq \Gamma_n$，所以每个
$!B \in \Gamma'$ 也 $\in \Gamma_n$，即 $\Gamma' \subseteq \Gamma_n$，而
$\Gamma_n$~一致。因此，每个有穷子集 $\Gamma' \subseteq \Gamma^*$ 都一致。
因
\iftag{FOL}{%
  \tagrefs{prfAX/{fol:axd:ptn:prop:proves-compact},
    prfSC/{fol:seq:ptn:prop:proves-compact},
    prfND/{fol:ntd:ptn:prop:proves-compact},
    prfTab/{fol:tab:ptn:prop:proves-compact}}}{%
  \tagrefs{prfAX/{pl:axd:ptn:prop:proves-compact},
    prfSC/{pl:seq:ptn:prop:proves-compact},
    prfND/{pl:ntd:ptn:prop:proves-compact},
    prfTab/{pl:tab:ptn:prop:proves-compact}}}，$\Gamma^*$ 一致。

$\Frm[L]$ 的每个!!{sentence}都出现在用来定义~$\Gamma^*$ 的列表上。若
$!A_n \notin \Gamma^*$，那是因为 $\Gamma_n \cup \{!A_n\}$ 不一致。但于是
$\lnot !A_n \in \Gamma^*$，所以 $\Gamma^*$ 是!!{complete}的。
\end{proof}

\end{document}
